\chapter{Boundary Values of Harmonic Functions}
\label{ch:harmonic}

A basic tool in the study of a function, $f$, on $E_n$ is obtained by the passage to a harmonic function, defined on the $(n+1)$-dimensional upper half-space, which has $f$ as its boundary value. In the first section we set down some of the useful basic facts concerning harmonic functions of several variables. The connection between an $f$, defined on $E_n$, and a harmonic function of $(n+1)$ variables (its Poisson integral) is studied in Section~2. Sections~3 and 4 deal with nontangential convergence at the boundary (for harmonic functions) and with harmonic majorization. While these notions are interesting in themselves, they will give us some important tools used in Chapters~III and VI.

\section{Basic Properties of Harmonic Functions}

We assume the reader is familiar with the properties of harmonic functions that are usually presented in a course in complex variables. Many of these results extend to higher dimensions; however, we can no longer rely on the theory of analytic functions (as is often done in two dimensions) to establish these properties. In this section, therefore, we develop some of the basic theory of harmonic functions of $n$ variables.

A function, $u$, defined in a domain (an open connected subset of $E_n$) is said to be \emph{harmonic} if it is twice differentiable (i.e., its second partial derivatives exist and are continuous) and satisfies the Laplace equation $\Delta u = \sum_{k=1}^{n} \partial^2 u / \partial x_k^2 = 0$ at every point of the domain. From the first two of the following examples of harmonic functions we can easily deduce that the Abel means and Poisson integrals introduced in the first section of Chapter~I are harmonic functions:

\begin{enumerate}
\item[(1)] For $t \in E_n$ fixed, $e^{-2\pi|t|y} e^{2\pi i x \cdot t} = u(x, y) = u(x_1, x_2, \dots, x_n, y)$ is harmonic in $E_{n+1}$;
\item[(2)] $P(x, y) = c_n [y/(|x|^2 + y^2)^{(n+1)/2}]$ is also harmonic in $E_{n+1}^+ = \{(x, y) = (x_1, \dots, x_n, y) \in E_{n+1} : y > 0\}$ (this domain will often be called the upper half-space);\footnote{We denote the points of $E_{n+1}^+$ by the ordered pairs $(x, y)$, where $x \in E_n$ and $y$ is a positive real, in order to emphasize the fact that this domain is a generalization of the upper half-plane (the case $n = 1$).}
\item[(3)] If $n \geq 3$ then $u(x) = |x|^{2-n}$ is harmonic in any domain not containing the origin. The corresponding example in two dimensions is $u(x) = \log |x|$. In a sense that will become clearer later, these functions ``generate'' all harmonic functions;
\item[(4)] We shall show later (Theorem~\ref{thm:c2-1-7}) that a harmonic function is infinitely differentiable, and, in fact, is real analytic (see (5.5) Chapter~IV). It is clear that any partial derivative of a harmonic function is also harmonic. As an illustration we observe that $P(x, y)$ is the partial derivative with respect to $y$ of $[c_n/(1 - n)] \cdot [1/(|x|^2 + y^2)^{(n-1)/2}]$, when $n > 1$, and of $c_n \log(|x|^2 + y^2)^{1/2}$, when $n$ is $1$;
\item[(5)] If $u$ is harmonic in a domain $\mathcal{D}$ then $\tau_h u$ is harmonic in $\mathcal{D} + h = \{x + h : x \in \mathcal{D}\}$. That is, harmonicity is preserved under translations. Similarly, harmonicity is preserved under rotations and dilations. (This can be computed directly, but will be more evident after we establish the characterization of harmonic functions in terms of the mean value property---see Theorem~\ref{thm:c2-1-1} below);
\item[(6)] $\sin(\sqrt{n} y) \prod_{k=1}^{n} \cosh x_k = u(x, y)$ is harmonic in $E_{n+1} = \{(x, y) : x \in E_n, y \in E_1\}$.
\end{enumerate}

We shall denote the mean value of a function $u$ taken over the surface of a sphere of radius $r > 0$ about the point $x$ by the symbols $\mathcal{M}(r) = \mathcal{M}_{x,u}(r)$ (depending on the necessity of showing its dependence on $u$, $x$ and $r$). That is, letting $\Sigma$ denote the surface of the unit sphere in $E_n$ and $dt'$ the element of surface area on $\Sigma$ (this was the notation used in the proof of Lemma~1.17 in Chapter~I)
\[
\mathcal{M}_{x,u}(r) = \frac{1}{\omega_{n-1}} \int_{\Sigma} u(x + r t')\,dt'.
\]

\begin{theorem}[Mean-Value Theorem for Harmonic Functions]\label{thm:c2-1-1}
If $u$ is harmonic in a domain $\mathcal{D}$ and if the sphere of radius $r_0$ about $x \in \mathcal{D}$ is contained in $\mathcal{D}$ then, for $0 < r \leq r_0$,
\[
u(x) = \mathcal{M}_{x,u}(r).
\]
\end{theorem}

\begin{proof}
Because of (5) we see that we can assume $x = 0$. Let $\mathcal{E}$ be a subdomain with a boundary $\partial\mathcal{E} \subset \mathcal{D}$ that is sufficiently smooth. Then, if we apply Green's theorem to the functions $u$ and $1$ in $\mathcal{E}$, we obtain
\begin{equation}\label{eq:c2-1-2}
\int_{\partial\mathcal{E}} \frac{\partial u}{\partial n}\,ds = 0,
\end{equation}
where $\partial/\partial n$ is differentiation in the direction of the outward directed normal to $\partial\mathcal{E}$ and $ds$ is the element of surface area on $\partial\mathcal{E}$. Let $\Sigma_\varepsilon$ and $\Sigma_r$, $0 < \varepsilon < r \leq r_0$, be the surfaces of the spheres of radii $\varepsilon$ and $r$ about the origin. By applying Green's second identity to the functions $u$ and
\[
v(x) = \begin{cases}
|x|^{-(n-2)} & \text{if } n > 2, \\
\log |x| & \text{if } n = 2
\end{cases}
\]
in the shell domain between these two spheres we obtain\footnote{The proof will be written out for the case $n > 2$. Obvious modifications give us the result when $n = 2$. When $n = 1$ the result is trivial.}
\[
0 = \left(\int_{\Sigma_r} - \int_{\Sigma_\varepsilon}\right) (-(n-2)|x|^{-(n-1)} u)\,ds - \left\{ \int_{\Sigma_r} v \frac{\partial u}{\partial n}\,ds - \int_{\Sigma_\varepsilon} v \frac{\partial u}{\partial n}\,ds \right\}.
\]
But $v$ is constant on $\Sigma_r$ (where it equals $r^{-(n-2)}$) and on $\Sigma_\varepsilon$ (where it equals $\varepsilon^{-(n-2)}$). Thus, by \eqref{eq:c2-1-2}, the term in the braces is $0$, and we obtain
\[
\varepsilon^{1-n} \int_{\Sigma_\varepsilon} u\,ds = r^{1-n} \int_{\Sigma_r} u\,ds.
\]
Hence,
\[
\mathcal{M}_{0,u}(r) = \frac{1}{\omega_{n-1}} \int_{\Sigma} u(rt')\,dt' = \frac{1}{\omega_{n-1}} \varepsilon^{1-n} \int_{\Sigma_\varepsilon} u\,ds.
\]
But the last expression obviously tends to $u(0)$ as $\varepsilon \to 0$; this proves the theorem.
\end{proof}

\begin{corollary}[Maximum Principle for Harmonic Functions]\label{cor:c2-1-3}
Suppose a real-valued harmonic function $u$, defined in a domain $\mathcal{D}$, attains its maximum $A$ in $\mathcal{D}$. Then $u$ is constant in $\mathcal{D}$.
\end{corollary}

\begin{proof}
Suppose $u(x) = A$ for some $x \in \mathcal{D}$. Then, by Theorem~\ref{thm:c2-1-1}, there exists $r_0 > 0$ such that $u(x) = \mathcal{M}_{x,u}(r)$ for all $0 < r \leq r_0$. Since $u \leq A$ throughout $\mathcal{D}$, we must have $u(x + rt') = A$ for all $t' \in \Sigma$ and $0 < r \leq r_0$. Thus, the set $\{y \in \mathcal{D} : u(y) = A\}$ is open in $\mathcal{D}$. Since $u$ is continuous, this set is also relatively closed in $\mathcal{D}$. Since $\mathcal{D}$ is connected, $u(y) = A$ for all $y \in \mathcal{D}$.
\end{proof}

By applying this result to $-u$ we obtain the \emph{minimum principle} for harmonic functions: Under the hypotheses of Corollary~\ref{cor:c2-1-3}, if $u$ attains its minimum in $\mathcal{D}$, then $u$ is constant in $\mathcal{D}$.

An equivalent form of the maximum (minimum) principle is the following:

\begin{customcor}{1.3$'$}\label{cor:c2-1-3p}
If $u$ is continuous on the closure $\overline{\mathcal{D}}$ of the bounded domain $\mathcal{D}$ and is harmonic in $\mathcal{D}$, then $u$ attains its maximum (and minimum) on the boundary $\partial\mathcal{D}$.
\end{customcor}

Applying this to the difference $u_1 - u_2$ we obtain:

\begin{corollary}\label{cor:c2-1-4}
If $u_1$ and $u_2$ are continuous on the closure $\overline{\mathcal{D}}$ of the bounded domain $\mathcal{D}$, are harmonic in $\mathcal{D}$, and $u_1 = u_2$ on the boundary $\partial\mathcal{D}$, then $u_1 = u_2$ throughout $\mathcal{D}$.
\end{corollary}

Another consequence of the mean value property is the following result, known as \emph{Liouville's theorem}:

\begin{theorem}\label{thm:c2-1-5}
If $v$ is harmonic on $E_n$ and bounded, then $v$ is a constant function.
\end{theorem}

\begin{proof}
Let $\Omega_n = \omega_{n-1}/n$ be the volume of the solid unit sphere of $E_n$. From the mean value property we have
\[
v(x) = v(x) \frac{n}{t^n} \int_0^t \rho^{n-1}\,d\rho = \frac{n}{t^n} \int_0^t \mathcal{M}_{x,v}(\rho) \rho^{n-1}\,d\rho = \frac{1}{t^n \Omega_n} \int_{|y-x| < t} v(y)\,dy.
\]
Thus, if $S_t(x)$ denotes the solid sphere about $x$ of radius $t$,
\[
|v(x_1) - v(x_2)| = \frac{1}{t^n \Omega_n} \left| \int_{S_t(x_1)} v(x)\,dx - \int_{S_t(x_2)} v(x)\,dx \right| \leq \frac{M}{t^n \Omega_n} \int_{(S_t(x_1) \Delta S_t(x_2))} dx,
\]
where $M = \sup_{x \in E_n} |v(x)|$ and $\Delta$ denotes the symmetric difference $(S_t(x_1) - S_t(x_2)) \cup (S_t(x_2) - S_t(x_1))$. But the measure of this symmetric difference divided by $t^n$ clearly tends to $0$ as $t \to \infty$, and it follows that $v(x_1) - v(x_2) = 0$.
\end{proof}

The mean value property completely characterizes harmonic functions. To see this we first establish the following lemma:

\begin{lemma}\label{lem:c2-1-6}
Suppose $u$ is twice differentiable in the domain $\mathcal{D}$ and $\mathcal{M}_{x,u}(r) = u(x)$ whenever the sphere of radius $r$ about $x$ is contained in $\mathcal{D}$. Then
\[
\mathcal{M}''(0) = \frac{1}{n} (\Delta u)(x),
\]
where $\mathcal{M}''(0) = \lim_{r \to 0} \mathcal{M}''(r)$. In particular, $u$ is harmonic in $\mathcal{D}$.
\end{lemma}

\begin{proof}
Fix $x \in \mathcal{D}$ and let $t' = (t_1', t_2', \dots, t_n') \in \Sigma$. We have
\[
\frac{d^2}{dr^2} \mathcal{M}(r) = \frac{1}{\omega_{n-1}} \int_{\Sigma} \left( \sum_{j,k=1}^n u_{jk}(x + rt') t_j' t_k' \right)\,dt',
\]
where $u_{jk} = \partial^2 u / \partial x_j \partial x_k$. Hence,
\[
\mathcal{M}''(0) = \frac{1}{\omega_{n-1}} \sum_{j,k=1}^n \int_{\Sigma} u_{jk}(x) t_j' t_k'\,dt'.
\]
But, by an easy symmetry argument, $\int_{\Sigma} t_j' t_k'\,dt' = 0$ if $j \neq k$, while $\int_{\Sigma} (t_j')^2\,dt' = (1/n) \int_{\Sigma} |t'|^2\,dt' = \omega_{n-1}/n$. Thus $\mathcal{M}''(0) = (1/n) \sum_{j=1}^n u_{jj}(x) = (1/n)(\Delta u)(x)$. Since $\mathcal{M}(r)$ is constant, $\mathcal{M}''(0) = 0$, which proves that $(\Delta u)(x) = 0$.
\end{proof}

The following result shows, in particular, that Lemma~\ref{lem:c2-1-6} is valid even if we weaken the hypothesis of second differentiability:

\begin{theorem}\label{thm:c2-1-7}
Suppose a continuous function, $u$, on a domain $\mathcal{D} \subset E_n$ satisfies the mean value property
\[
u(x) = \mathcal{M}_{x,u}(r)
\]
whenever the sphere of radius $r$ about $x$ is contained in $\mathcal{D}$. Then $u$ is infinitely differentiable, and is harmonic in $\mathcal{D}$.
\end{theorem}

\begin{proof}
Since the problem is local we can assume, by restricting $u$ to a sphere whose closure belongs to $\mathcal{D}$, that $u$ is bounded and continuous on $E_n$, and that the mean value property holds for all $x \in E_n$ and all $r \leq r_0$. Let $\psi \in C_0^\infty(E_n)$ be a non-negative radial function supported in $\{x : |x| \leq r_0\}$ with $\int_{E_n} \psi(x)\,dx = 1$. Then $\psi(x) = \psi_0(|x|)$. Using polar coordinates,
\[
(u * \psi)(x) = \int_0^{r_0} \psi_0(r) \left( \int_{|y|=r} u(x - y)\,d\sigma(y) \right)\,dr = \omega_{n-1} \int_0^{r_0} \psi_0(r) r^{n-1} \mathcal{M}_{x,u}(r)\,dr = u(x) \int_{E_n} \psi(y)\,dy = u(x).
\]
Since $\psi \in C_0^\infty$, $u * \psi$ is infinitely differentiable, and thus $u$ is $C^\infty$. By Lemma~\ref{lem:c2-1-6}, $u$ is harmonic.
\end{proof}

\begin{corollary}\label{cor:c2-1-8}
If $\{u_k\}$ is a sequence of harmonic functions in $\mathcal{D}$ that converges uniformly on compact subsets of $\mathcal{D}$ to a function $u$, then $u$ is harmonic in $\mathcal{D}$.
\end{corollary}

\begin{proof}
Uniform convergence preserves continuity and commutes with the integral in the mean value formula. Hence $u$ is continuous and satisfies the mean value property. By Theorem~\ref{thm:c2-1-7}, $u$ is harmonic.
\end{proof}

Next we consider the Dirichlet problem for a ball: given a continuous function on the boundary of a ball, find a function harmonic in the interior that agrees with the given function on the boundary.

\begin{theorem}\label{thm:c2-1-9}
Let $B = \{x \in E_n : |x| < 1\}$ be the unit ball and $S = \partial B$. For $x \in B$ and $t' \in S$, define the Poisson kernel for the ball by
\[
P(x, t') = \frac{1 - |x|^2}{|x - t'|^n}.
\]
If $f$ is continuous on $S$, the function
\[
u(x) = \begin{cases}
\dfrac{1}{\omega_{n-1}} \displaystyle\int_S P(x, t') f(t')\,dt' & \text{for } |x| < 1, \\
f(x) & \text{for } |x| = 1
\end{cases}
\]
is continuous on $\overline{B}$ and harmonic in $B$.
\end{theorem}

\begin{theorem}\label{thm:c2-1-10}
The Poisson kernel $P(x, t')$ satisfies:
\begin{enumerate}
\item[(i)] $P(x, t') > 0$ for all $|x| < 1$ and $t' \in S$;
\item[(ii)] $\dfrac{1}{\omega_{n-1}} \displaystyle\int_S P(x, t')\,dt' = 1$ for all $|x| < 1$;
\item[(iii)] For every $\delta > 0$ and $t_0 \in S$, $\lim_{x \to t_0, |x|<1} P(x, t') = 0$ uniformly for $|t' - t_0| \geq \delta$.
\end{enumerate}
\end{theorem}

\begin{corollary}[Harnack's Inequality]\label{cor:c2-1-11}
If $u$ is non-negative and harmonic in the unit ball $B$, then for all $|x| < 1$,
\[
\frac{1 - |x|}{(1 + |x|)^{n-1}} u(0) \leq u(x) \leq \frac{1 + |x|}{(1 - |x|)^{n-1}} u(0).
\]
\end{corollary}

\begin{theorem}[Harnack's Theorem]\label{thm:c2-1-12}
Let $\{u_k\}$ be a monotonically increasing sequence of harmonic functions in a domain $\mathcal{D}$. Then either $u_k(x) \to \infty$ for all $x \in \mathcal{D}$, or $\{u_k\}$ converges uniformly on compact subsets of $\mathcal{D}$ to a harmonic function.
\end{theorem}

\begin{theorem}\label{thm:c2-1-13}
Let $u$ be harmonic in $E_{n+1}^+ = \{(x, y) : x \in E_n, y > 0\}$. If $u$ is bounded and continuous on $\overline{E_{n+1}^+}$, then
\begin{equation}\label{eq:c2-1-14}
u(x, y) = (P_y * f)(x) = \int_{E_n} P(x - t, y) f(t)\,dt,
\end{equation}
where $f(x) = u(x, 0)$ and $P(x, y) = c_n y / (|x|^2 + y^2)^{(n+1)/2}$.
\end{theorem}

\begin{corollary}\label{cor:c2-1-15}
If $u$ is bounded and harmonic in $E_{n+1}^+$ and $\lim_{y \to 0} u(x, y) = 0$ for all $x \in E_n$, then $u \equiv 0$.
\end{corollary}

\section{The Characterization of Poisson Integrals}

In this section we study the relationship between functions $f$ on $E_n$ and harmonic functions on the upper half-space $E_{n+1}^+$.

\begin{theorem}\label{thm:c2-2-1}
Let $1 \leq p \leq \infty$. A harmonic function $u(x, y)$ in $E_{n+1}^+$ is the Poisson integral of a function $f \in L^p(E_n)$,
\begin{equation}\label{eq:c2-2-2}
u(x, y) = \int_{E_n} P(x - t, y) f(t)\,dt,
\end{equation}
if and only if:
\begin{enumerate}
\item[(a)] For $1 < p \leq \infty$, $\sup_{y > 0} \|u(\,\cdot\,, y)\|_p < \infty$;
\item[(b)] For $p = 1$, $u(x, y) \to f(x)$ in $L^1(E_n)$ norm as $y \to 0$.
\end{enumerate}
In this case, $\|u(\,\cdot\,, y)\|_p \leq \|f\|_p$ for all $y > 0$, and $\lim_{y \to 0} \|u(\,\cdot\,, y) - f\|_p = 0$ for $1 \leq p < \infty$.
\end{theorem}

\begin{theorem}\label{thm:c2-2-3}
A harmonic function $u(x, y)$ in $E_{n+1}^+$ is the Poisson--Stieltjes integral of a finite Borel measure $\mu$ on $E_n$,
\begin{equation}\label{eq:c2-2-4}
u(x, y) = \int_{E_n} P(x - t, y)\,d\mu(t),
\end{equation}
if and only if $\sup_{y > 0} \|u(\,\cdot\,, y)\|_1 < \infty$.
\end{theorem}

\begin{theorem}\label{thm:c2-2-5}
A harmonic function $u(x, y) \geq 0$ on $E_{n+1}^+$ is the Poisson--Stieltjes integral of a unique non-negative Borel measure $\mu$ on $E_n$ plus a linear term $cy$, $c \geq 0$.
\end{theorem}

\begin{lemma}\label{lem:c2-2-6}
Let $\{\mu_k\}$ be a sequence of finite Borel measures on $E_n$ such that $\sup_k \|\mu_k\| < \infty$. Then there exists a subsequence $\{\mu_{k_j}\}$ and a finite measure $\mu$ such that $\mu_{k_j} \to \mu$ weakly; that is,
\[
\lim_{j \to \infty} \int_{E_n} g(t)\,d\mu_{k_j}(t) = \int_{E_n} g(t)\,d\mu(t)
\]
for every continuous function $g$ vanishing at infinity.
\end{lemma}

\begin{lemma}\label{lem:c2-2-7}
If $u$ is harmonic in $E_{n+1}^+$ and $\sup_{y > 0} \|u(\,\cdot\,, y)\|_p < \infty$ for some $1 \leq p \leq \infty$, then for each fixed $y_0 > 0$, $u(x, y + y_0)$ is the Poisson integral of the function $u(x, y_0)$. That is, the semigroup property
\[
u(x, y_1 + y_2) = \int_{E_n} P(x - t, y_2) u(t, y_1)\,dt
\]
holds for all $y_1, y_2 > 0$.
\end{lemma}

\section{The Hardy--Littlewood Maximal Function and Nontangential Convergence of Harmonic Functions}

Suppose $f$ is a measurable function on $E_n$. For each $x \in E_n$ we consider the mean value
\[
\frac{1}{|S_x|} \int_{S_x} |f(t)|\,dt, \qquad \frac{1}{|Q_x|} \int_{Q_x} |f(t)|\,dt,
\]
where $S_x$ ranges over all spheres with center $x$ and $Q_x$ over all cubes with center $x$ and sides parallel to the coordinate axes. We define the \emph{spherical maximal function} $m_f$ and the \emph{cubic maximal function} $\overline{m}_f$ by
\begin{equation}\label{eq:c2-3-1}
m_f(x) = \sup_{S_x} \frac{1}{|S_x|} \int_{S_x} |f(t)|\,dt, \qquad
\overline{m}_f(x) = \sup_{Q_x} \frac{1}{|Q_x|} \int_{Q_x} |f(t)|\,dt.
\end{equation}
It is clear that, if $f$ is essentially bounded, then $m_f(x)$ and $\overline{m}_f(x)$ do not exceed $\|f\|_\infty$. We shall show that $m_f(x)$ and $\overline{m}_f(x)$ are finite almost everywhere whenever $f \in L^p(E_n)$, $1 \leq p \leq \infty$.

Given a sphere $S_x$ of radius $r > 0$ and center $x$, let $q_x$ be the inscribed cube (it has, therefore, diagonals of length $2r$) and $Q_x$ the circumscribed cube (it has sides of length $2r$), both cubes having sides parallel to the axes. Then, clearly, there exist constants $a_n$ and $A_n$, depending only on the dimension, such that $|Q_x| \leq A_n |S_x|$ and $|S_x| \leq a_n |q_x|$. Thus,
\[
\frac{1}{|S_x|} \int_{S_x} |f(t)|\,dt \leq \frac{A_n}{|Q_x|} \int_{Q_x} |f(t)|\,dt.
\]
But this implies
\begin{equation}\label{eq:c2-3-2}
m_f(x) \leq A_n \overline{m}_f(x).
\end{equation}
Similarly, we have
\begin{equation}\tag{3.2$'$}\label{eq:c2-3-2p}
\overline{m}_f(x) \leq a_n m_f(x).
\end{equation}
We shall be more interested in the spherical maximal function $m_f$. In some cases, however, it is easier to establish needed estimates for $\overline{m}_f$; inequalities \eqref{eq:c2-3-2} and \eqref{eq:c2-3-2p} can then be used to obtain corresponding results for $m_f$. The following lemma will enable us to show that $\overline{m}_f(x) < \infty$ almost everywhere (and, hence, $m_f(x) < \infty$ a.e.).

\begin{lemma}\label{lem:c2-3-3}
Let $R_1 \subset R_2 \subset \dots \subset R_k$ be nonempty open rectangles in $E_n$ with sides parallel to the axes and centered at $0$. Suppose $S \subset E_n$ is a bounded set and, for each $x$ in $S$, let $i(x)$ be an integer satisfying $1 \leq i(x) \leq k$. If
\[
R^x = \{u \in E_n : u - x \in R_{i(x)}\},
\]
then there exist a finite number of points, $x_1, x_2, \dots, x_l$, of $S$ such that $S \subset R^{x_1} \cup \cdots \cup R^{x_l}$, and each $v$ in $E_n$ belongs to at most $2^n$ of the sets $R^{x_1}, \dots, R^{x_l}$.
\end{lemma}

\begin{proof}
Choose $x_1$ so that $i(x_1)$ is the largest possible (we can do this since $i(x) \leq k$). Having done so, we then choose $x_2 \in S - R^{x_1}$ so that $i(x_2)$ is largest possible and continue in this way by choosing $x_3 \in S - \{R^{x_1} \cup R^{x_2}\}$ so that $i(x_3)$ is maximal, etc. We thus obtain a sequence of open rectangles $R^{x_1}, R^{x_2}, \dots$ centered at $x_1, x_2, \dots$ with the property that each center $x_j$ belongs to no other rectangle $R^{x_i}$, $i \neq j$. But this means that the distance between $x_i$ and $x_j$, $i \neq j$, must exceed half the length of the shortest side of $R_1$ (the smallest rectangle). Since, by assumption, this half-length is positive and the set $S$ is bounded it follows that the sequence $x_1, x_2, \dots$ is finite and, after reaching the, say, $l$-th stage of our process we have obtained a covering $\{R^{x_1}, \dots, R^{x_l}\}$ of $S$.

It now remains to be shown that a point $v \in E_n$ can belong to no more than $2^n$ of the rectangles $R^{x_1}, \dots, R^{x_l}$. In order to see this, consider the hyperplanes through $v$ that are parallel to the coordinate hyperplanes. These form $2^n$ ``octants'' with vertex $v$ that cover all of $E_n$. Suppose two of the points, $x_i$ and $x_j$ ($1 \leq i, j \leq l$), belong to the same octant in such a way that $R^{x_i}$ and $R^{x_j}$ both contain $v$. Since each of these two rectangles is a translate of one of the rectangles $R_1, \dots, R_k$, one of these two, say $R^{x_i}$, must have sides that are at least as long as those of the other, $R^{x_j}$. But this, together with the fact that $v \in R^{x_i} \cap R^{x_j}$, implies that $x_j \in R^{x_i}$, contradicting the fact that each center $x_j$ belongs to no other rectangle $R^{x_i}$. Thus, at most one rectangle with center in any given ``octant'' can contain $v$ and the lemma is proved.
\end{proof}

\begin{theorem}\label{thm:c2-3-4}
There exists a constant $c = c(n)$, depending only on the dimension, such that whenever $f \in L^1(E_n)$ and
\[
F_s = \{x \in E_n : m_f(x) > s > 0\}
\]
then
\[
|F_s| \leq \frac{c\|f\|_1}{s},
\]
where $|F_s|$ denotes the Lebesgue measure of $F_s$. In particular, $m_f(x) < \infty$ for almost every $x \in E_n$.
\end{theorem}

\begin{proof}
We shall show that
\begin{equation}\label{eq:c2-3-5}
\left|\{x \in E_n : \overline{m}_f(x) > s\}\right| \leq 2^n \frac{\|f\|_1}{s}
\end{equation}
for each $s > 0$. Since, by \eqref{eq:c2-3-2}, $F_s \subset \{x \in E_n : \overline{m}_f(x) > s/A_n\}$, we would then have
\[
|F_s| \leq \left|\left\{x \in E_n : \overline{m}_f(x) > \frac{s}{A_n}\right\}\right| \leq \frac{2^n A_n \|f\|_1}{s}
\]
and the theorem follows with $c = 2^n A_n$.

Let $S$ be any compact subset of $\{x \in E_n : \overline{m}_f(x) > s\}$. For each $x \in S$, therefore, there exists a cube $Q_x$ such that
\[
\frac{1}{|Q_x|} \int_{Q_x} |f(t)|\,dt > s.
\]
By continuity of the integral, there exists a neighborhood $U$ of $x$ such that for each $u \in U$
\[
\frac{1}{|Q_u|} \int_{Q_u} |f(t)|\,dt > s,
\]
where $Q_u = \{(w - x) + u : w \in Q_x\}$ is the cube of the same size as $Q_x$ that is centered at $u$. Since $S$ is compact there are a finite number of such neighborhoods covering $S$. Hence, there are a finite number of cubes $Q_1 \subset \cdots \subset Q_k$, centered at $0$, such that for each $x \in S$ there exists an integer $i(x)$, $1 \leq i(x) \leq k$, for which
\[
\frac{1}{|Q^x|} \int_{Q^x} |f(t)|\,dt > s,
\]
where $Q^x = \{u \in E_n : u - x \in Q_{i(x)}\}$. By Lemma~\ref{lem:c2-3-3}, we can find a finite number of points $x_1, \dots, x_l$ in $S$ such that
\[
S \subset Q^{x_1} \cup \cdots \cup Q^{x_l}
\]
and each $v \in E_n$ belongs to at most $2^n$ of these cubes. If, for a set $A$, $\chi_A$ denotes the characteristic function of $A$, this last property is equivalent to the inequality
\[
\sum_{j=1}^l \chi_{Q^{x_j}} \leq 2^n \chi_{\bigcup_{j=1}^l Q^{x_j}}.
\]
Thus,
\[
\begin{aligned}
|S| &\leq \left|\bigcup_{j=1}^l Q^{x_j}\right| \leq \sum_{j=1}^l |Q^{x_j}| \leq \frac{1}{s} \sum_{j=1}^l \int_{Q^{x_j}} |f(t)|\,dt \\
&= \frac{1}{s} \int_{E_n} \left\{\sum_{j=1}^l \chi_{Q^{x_j}}(t)\right\} |f(t)|\,dt \leq \frac{2^n}{s} \int_{E_n} \chi_{\bigcup_{j=1}^l Q^{x_j}}(t) |f(t)|\,dt \leq \frac{2^n}{s} \|f\|_1.
\end{aligned}
\]
Since $S$ is an arbitrary compact subset of $\{x \in E_n : \overline{m}_f(x) > s\}$ inequality \eqref{eq:c2-3-5} and, hence, the theorem now follow immediately.
\end{proof}

We have already observed that $m_g(x) \leq \|g\|_\infty < \infty$ for any $g \in L^\infty(E_n)$. From Theorem~\ref{thm:c2-3-4} we have $m_f(x) < \infty$ for almost every $x$ when $f \in L^1(E_n)$. It follows from these two facts that $m_h(x) < \infty$ a.e. whenever $h \in L^p(E_n)$, $1 < p < \infty$. In order to see this we first observe that the maximal operator is an example of a \emph{sublinear} operator. By this we mean an operator $T$ mapping a linear space of measurable functions defined on a measure space into measurable functions defined on (possibly) another measure space satisfying
\[
|[T(f+g)](x)| \leq |(Tf)(x)| + |(Tg)(x)|
\]
almost everywhere (this property is called subadditivity), and
\[
|T(af)| = |a| |Tf|
\]
whenever $a$ is a scalar and $f$ is in the domain of $T$. Thus, if, for $h \in L^p(E_n)$, we put
\[
g(x) = \begin{cases} h(x) & \text{when } |h(x)| \leq 1, \\ 0 & \text{when } |h(x)| > 1, \end{cases}
\]
and $f = h - g$, then $g \in L^\infty(E_n)$ and $f \in L^1(E_n)$. By property (i) (subadditivity), therefore,
\[
m_h(x) = m_{f+g}(x) \leq m_f(x) + m_g(x).
\]
Consequently, $m_h(x) < \infty$ for almost every $x$. All these observations are, of course, equally valid for the cubic maximal function.

The maximal operator, as a mapping defined on $L^1(E_n)$, is not a bounded transformation into $L^1(E_n)$ (e.g., if $f$ is the characteristic function of the interval $[0, 1]$ then $m_f$ is not even integrable). Theorem~\ref{thm:c2-3-4}, however, tells us that it is ``almost'' bounded. In order to make this statement more precise we introduce the notion of the \emph{distribution function} $\lambda = \lambda_g$ of a function $g \in L^p(E_n)$: if $s > 0$ then $\lambda(s)$ is the Lebesgue measure of the set $F_s = \{x \in E_n : |g(x)| > s\}$. We thus obtain a nonincreasing function $\lambda$ on the positive real axis. It is clear that $s^p \lambda(s) = \int_{F_s} s^p\,dx \leq \int_{F_s} |g|^p\,dx \leq \|g\|_p^p$. That is, $\lambda(s)$ is majorized by a constant times $s^{-p}$. This condition is clearly not sufficient to insure the integrability of $|g|^p$. On the other hand, the equality
\begin{equation}\label{eq:c2-3-6}
\|g\|_p^p = -\int_0^\infty s^p\,d\lambda(s) = p \int_0^\infty s^{p-1} \lambda(s)\,ds,
\end{equation}
which follows from integration by parts, shows that the integrability of $|g|^p$ is equivalent to the convergence of $\int_0^\infty s^{p-1}\lambda(s)\,ds$.

\begin{theorem}\label{thm:c2-3-7}
There exists a constant $b = b(p, n)$, depending on the dimension $n$ and the exponent $p > 1$, such that
\[
\|m_f\|_p \leq b\|f\|_p
\]
for all $f \in L^p(E_n)$.
\end{theorem}

\begin{proof}
Suppose $f \in L^p(E_n)$, $1 < p \leq \infty$, and $s > 0$. For each $x \in E_n$, let
\[
f^s(x) = \begin{cases} f(x) & \text{if } |f(x)| > s, \\ 0 & \text{if } |f(x)| \leq s, \end{cases}
\]
and
\[
f_s(x) = \begin{cases} 0 & \text{if } |f(x)| > s, \\ f(x) & \text{if } |f(x)| \leq s. \end{cases}
\]
Then $f = f^s + f_s$, with $f^s \in L^1(E_n)$ and $f_s \in L^\infty(E_n)$.

We let $\lambda$, $\lambda^s$, $\lambda_s$ be the distribution functions of $m_f$, $m_{f^s}$, and $m_{f_s}$. Since, by subadditivity, $m_f \leq m_{f^s} + m_{f_s}$ it follows that $\lambda(2s) \leq \lambda^s(s) + \lambda_s(s)$. We already observed that the maximal function of an essentially bounded function $g$ is bounded by $\|g\|_\infty$. Thus, $m_{f_s}(x) \leq \|f_s\|_\infty \leq s$ and it follows that $\lambda_s(s) = 0$. Therefore $\lambda(2s) \leq \lambda^s(s)$.

From \eqref{eq:c2-3-6}, this last inequality and Theorem~\ref{thm:c2-3-4} we therefore obtain
\[
\begin{aligned}
\|m_f\|_p^p &= p2^p \int_0^\infty s^{p-1} \lambda(2s)\,ds \leq p2^p \int_0^\infty s^{p-1} \lambda^s(s)\,ds \\
&\leq p2^p \int_0^\infty s^{p-1} c s^{-1} \left\{\int_{E_n} |f^s(x)|\,dx\right\} ds
= p2^p c \int_0^\infty s^{p-2} \left\{\int_{|f(x)| > s} |f(x)|\,dx\right\} ds \\
&= p2^p c \int_{E_n} |f(x)| \left\{\int_0^{|f(x)|} s^{p-2}\,ds\right\} dx
= \frac{p2^p c}{p-1} \int_{E_n} |f(x)| |f(x)|^{p-1}\,dx,
\end{aligned}
\]
which proves the theorem with $b = b(p, n) = 2[pc/(p-1)]^{1/p}$.
\end{proof}

Another way of expressing $m_f$ is the following one. Let $\varphi$ be the characteristic function of the solid unit sphere $\{x \in E_n : |x| \leq 1\}$ divided by $\Omega_n$ (so that $\int_{E_n} \varphi(x)\,dx = 1$). Let $\varepsilon > 0$ and put $\varphi_\varepsilon(t) = \varepsilon^{-n}\varphi(t/\varepsilon)$. Then, if $f \in L^p(E_n)$, $1 \leq p \leq \infty$, and if, for the sake of notational simplicity, we assume $f \geq 0$, we have
\[
(f * \varphi_\varepsilon)(x) = \frac{1}{\varepsilon^n \Omega_n} \int_{|t| \leq \varepsilon} f(x - t)\,dt;
\]
therefore,
\begin{equation}\label{eq:c2-3-8}
m_f(x) = \sup_{\varepsilon > 0} (f * \varphi_\varepsilon)(x).
\end{equation}
Similarly, if $\varphi$ is the characteristic function of the cube $\{x = (x_1, \dots, x_n) \in E_n : -\frac{1}{2} \leq x_j \leq \frac{1}{2}, j = 1, 2, \dots, n\}$, then the right side of \eqref{eq:c2-3-8} gives us $\overline{m}_f(x)$. We have seen that $\overline{m}_f$ and $m_f$ are ``equivalent'' in the sense that inequalities \eqref{eq:c2-3-2} and \eqref{eq:c2-3-2p} hold. It is natural, therefore, to ask about the relation of $m_f$ with the supremum in \eqref{eq:c2-3-8} when $\varphi \geq 0$ is a more general function.

For example, suppose $\varphi = \sum_{k=1}^m c_k \chi_k$, $c_k > 0$, where $\chi_k$ is the characteristic function of the sphere $\sigma_k = \{x \in E_n : |x| \leq r_k\}$. Then, for $f \geq 0$,
\[
\begin{aligned}
(f * \varphi_\varepsilon)(x) &= \sum_{k=1}^m c_k \varepsilon^{-n} \int_{|t| \leq \varepsilon r_k} f(x - t)\,dt \\
&= \sum_{k=1}^m c_k r_k^n \Omega_n \frac{1}{\Omega_n (\varepsilon r_k)^n} \int_{|t| \leq \varepsilon r_k} f(x - t)\,dt
\leq m_f(x) \sum_{k=1}^m c_k |\sigma_k| = m_f(x) \|\varphi\|_1.
\end{aligned}
\]
That is,
\begin{equation}\label{eq:c2-3-9}
\sup_{\varepsilon > 0} (f * \varphi_\varepsilon)(x) \leq m_f(x) \|\varphi\|_1.
\end{equation}
This result clearly extends to any $\varphi \in L^1(E_n)$ of the form $\varphi(t) = \psi(|t|)$ (i.e., $\varphi$ is radial) with $\psi$ a nonnegative decreasing function on $[0, \infty)$. We need only approximate $\varphi$ from below by an ascending sequence of simple functions of the type just considered and apply the Lebesgue monotone convergence theorem.

In particular, if we choose $\varphi(t) = c_n/(|t|^2 + 1)^{(n+1)/2}$ then $u(x, y) = (f * \varphi_y)(x)$, $y > 0$, is the Poisson integral of $f$ and we have the theorem:

\begin{theorem}\label{thm:c2-3-10}
If $u(x, y)$, $y > 0$, is the Poisson integral of $f \in L^p(E_n)$, $1 \leq p \leq \infty$, then $|u(x, y)| \leq m_f(x)$.
\end{theorem}

It is not hard to establish the converse to this last result:

\begin{theorem}\label{thm:c2-3-11}
If $f \geq 0$ and its Poisson integral
\[
u(x, y) = \int_{E_n} f(x - t) P(t, y)\,dt, \qquad y > 0,
\]
is defined (for example, if $f$ belongs to $L^p(E_n)$, $1 \leq p \leq \infty$), then there exists a constant $A = A(n)$, depending only on the dimension, such that
\[
m_f(x) \leq A \left\{\sup_{y > 0} u(x, y)\right\}.
\]
\end{theorem}

\begin{proof}
Fix $r > 0$, then
\[
\begin{aligned}
\sup_{y > 0} u(x, y) \geq u(x, r) &= c_n \int_{E_n} f(x - t) \frac{r}{(|t|^2 + r^2)^{(n+1)/2}}\,dt
\geq c_n' \int_{|t| \leq r} f(x - t) r^{-n}\,dt \\
&= c_n' \Omega_n \left\{\frac{1}{\Omega_n r^n} \int_{|t| \leq r} f(x - t)\,dt\right\}.
\end{aligned}
\]
Thus, taking the supremum over all $r > 0$ of the expressions on the right we obtain the theorem with $A(n) = 1/(c_n' \Omega_n)$.
\end{proof}

The importance of the maximal function operator lies in the fact that it majorizes many important operators in analysis. We have just seen an example of this situation in the case of the Poisson integral. The following theorem gives us a method for using this majorization property in order to obtain pointwise convergence results for many of the families of operators that we shall encounter.

\begin{theorem}\label{thm:c2-3-12}
Let $\{T_\varepsilon\}$, $0 < \varepsilon$, be a family of linear operators mapping $L^p(E_n)$, $1 \leq p \leq \infty$, into the space of measurable functions on $E_n$. For each $h \in L^p(E_n)$ define $Mh$ by letting $(Mh)(x) = \sup_{\varepsilon > 0} |(T_\varepsilon h)(x)|$ for each $x \in E_n$.

Suppose that there exists a constant $a > 0$ and a real number $q \geq 1$ such that
\begin{equation}\label{eq:c2-3-13}
\left|\left\{x : (Mh)(x) > \lambda\right\}\right| \leq (a\|h\|_p \lambda^{-1})^q
\end{equation}
for all $\lambda > 0$ and $h \in L^p(E_n)$. If there exists a dense subset $\mathcal{D}$ of $L^p(E_n)$ such that $\lim_{\varepsilon \to 0} (T_\varepsilon g)(x)$ exists and is finite almost everywhere whenever $g \in \mathcal{D}$, then for each $f \in L^p(E_n)$, $\lim_{\varepsilon \to 0} (T_\varepsilon f)(x)$ exists and is finite almost everywhere.
\end{theorem}

\begin{proof}
Let $f$ be a function in $L^p(E_n)$. For each positive integer $k$, let $F_k$ be the set of all $x \in E_n$ such that $|(T_{\varepsilon'} f)(x) - (T_{\varepsilon''} f)(x)| > 2k^{-1}$ for infinitely many pairs $(\varepsilon', \varepsilon'')$ tending to $(0, 0)$. For each $\eta > 0$ we can find $g$ and $h$ in $L^p(E_n)$ such that $f = g + h$ with $g \in \mathcal{D}$ and $\|h\|_p < \eta$. We shall show that $|F_k| \leq (2a(k+1)\eta)^q$; since $\eta$ is arbitrary this means that $F_k$ has measure $0$. Thus, $F = \bigcup_{k=1}^\infty F_k$ has measure $0$ and, for $x \in E_n - F$, $\lim_{\varepsilon \to 0} (T_\varepsilon f)(x)$ exists and is finite. This, then, would prove the theorem.

Let $G$ be the set of all $x \in E_n$ for which $\lim_{\varepsilon \to 0} (T_\varepsilon g)(x)$ exists and is finite. Since $E_n - G$ has measure zero it suffices to show that $|F_k \cap G| \leq (2a(k+1)\eta)^q$. Since
\[
(T_{\varepsilon'} f)(x) - (T_{\varepsilon''} f)(x) = \{(T_{\varepsilon'} h)(x) - (T_{\varepsilon''} h)(x)\} + \{(T_{\varepsilon'} g)(x) - (T_{\varepsilon''} g)(x)\}
\]
and, for $x \in G$, $\lim_{\varepsilon', \varepsilon'' \to 0} \{(T_{\varepsilon'} g)(x) - (T_{\varepsilon''} g)(x)\} = 0$, it follows that $F_k \cap G$ is contained in the set of all $x \in E_n$ such that $|(T_{\varepsilon'} h)(x) - (T_{\varepsilon''} h)(x)| \geq k^{-1}$ for infinitely many pairs $(\varepsilon', \varepsilon'')$ tending to $(0,0)$. This set, in turn, is contained in $\{x : (Mh)(x) > (2k)^{-1}\}$, whose measure is at most $(2a k \eta)^q \leq (2a(k+1)\eta)^q$ by \eqref{eq:c2-3-13}.
\end{proof}

An example of such a family is given by defining $T_\varepsilon f$, for $\varepsilon > 0$, by letting
\[
(T_\varepsilon f)(x) = (\varepsilon^n \Omega_n)^{-1} \int_{|t| < \varepsilon} f(x - t)\,dt
\]
for each $x \in E_n$. In this case we have $Mf = m_f$ and, by Theorem~\ref{thm:c2-3-4}, we have
\[
\left|\left\{x : (Mf)(x) > \lambda\right\}\right| \leq c\|f\|_1 \lambda^{-1}
\]
for $f \in L^1(E_n)$. Furthermore, if $g$ belongs to the class of continuous functions with compact support (a dense subset of $L^1(E_n)$) then it is clear that $\lim_{\varepsilon \to 0} (T_\varepsilon g)(x) = g(x)$. Thus, by Theorem~\ref{thm:c2-3-12}, $\lim_{\varepsilon \to 0} (T_\varepsilon f)(x)$ exists and is finite for almost all $x$. From this we easily deduce that almost every $x$ is a point of differentiability of the integral of $f$ (see Chapter~I, (1.23) and (1.24)); more precisely, we have:

\begin{corollary}\label{cor:c2-3-14}
If $f$ is locally integrable on $E_n$ (i.e., $f$ is integrable over any bounded subset of $E_n$) then
\[
\lim_{\varepsilon \to 0} \varepsilon^{-n} \int_{|t| < \varepsilon} [f(x - t) - f(x)]\,dt = 0
\]
for almost every $x \in E_n$.
\end{corollary}

\begin{proof}
Since the problem is local we can assume $f \in L^1(E_n)$ (otherwise we can multiply $f$ by the characteristic function of a sphere of radius $r$ about $0$ and obtain the desired convergence at almost every point within this sphere; letting $r$ tend to $\infty$ we then obtain it for almost all points of $E_n$). We already know that
\[
\lim_{\varepsilon \to 0} (T_\varepsilon f)(x) = \lim_{\varepsilon \to 0} (\varepsilon^n \Omega_n)^{-1} \int_{|t| < \varepsilon} f(x - t)\,dt
\]
exists and is finite for almost all $x$. We only need to show that this limit is $f(x)$ a.e. But
\[
\begin{aligned}
\|T_\varepsilon f - f\|_1 &= \int_{E_n} \left|\frac{1}{\varepsilon^n \Omega_n} \int_{|t| < \varepsilon} [f(x - t) - f(x)]\,dt\right| dx \\
&\leq \frac{1}{\varepsilon^n \Omega_n} \int_{|t| < \varepsilon} \left\{\int_{E_n} |f(x - t) - f(x)|\,dx\right\} dt \to 0
\end{aligned}
\]
as $\varepsilon \to 0$ since the $L^1$ modulus of continuity $\omega_{1,f}(t) \to 0$ as $|t| \to 0$ (see the proof of Theorem~1.18 in Chapter~I). Consequently, there exists a sequence, $\{\varepsilon_k\}$, tending to $0$ such that $(T_{\varepsilon_k} f)(x) \to f(x)$ a.e. as $k \to \infty$. It follows, therefore, that $\lim_{\varepsilon \to 0} (T_\varepsilon f)(x) = f(x)$ a.e.
\end{proof}

This method extends immediately to Poisson integrals and we can obtain from it a direct proof of Theorem~\ref{thm:c2-2-1}. Our intention, however, is to extend this theorem to include the more general nontangential approach to the boundary that was mentioned briefly at the beginning of this section. We first make this notion more precise and then will show that Poisson integrals have these more general boundary values.

If $\alpha > 0$ and $x_0 \in E_n$ the \emph{cone} in $E_{n+1}^+$ with vertex $(x_0, 0)$ and aperture $\alpha$ is the region
\[
\Gamma_\alpha(x_0) = \{(x, y) \in E_{n+1}^+ : |x - x_0| < \alpha y\}.
\]
If $u$ is defined in $E_{n+1}^+$ we shall say that it has a \emph{nontangential limit} $l$ at $x_0 \in E_n$ if, for each $\alpha > 0$,
\begin{equation}\label{eq:c2-3-15}
\lim_{\substack{(x,y) \to (x_0,0)}} u(x, y) = l,
\end{equation}
where $(x, y)$ tends to $(x_0, 0)$ within the cone $\Gamma_\alpha(x_0)$.

The importance of the nontangential approach to the boundary will become clearer later. In several basic results that we shall develop, the vertical approach to the boundary is not sufficient (see, e.g., (5.4)).

\begin{theorem}\label{thm:c2-3-16}
If $f \in L^p(E_n)$, $1 \leq p \leq \infty$, then the Poisson integral
\[
\int_{E_n} f(x - t) P(t, y)\,dt = u(x, y)
\]
has the nontangential limit $f(x_0)$ at each $x_0$ belonging to the Lebesgue set of $f$; in particular, it has these limits at almost all points of $E_n$.
\end{theorem}

\begin{proof}
We first observe that the proof of Theorem~1.25 in the first chapter consisted in showing that
\[
0 = \lim_{y \to 0} \int_{E_n} |f(x_0 - t) - f(x_0)| |\varphi_y(t)|\,dt = \lim_{y \to 0} \int_{E_n} |f(t) - f(x_0)| |\varphi_y(x_0 - t)|\,dt
\]
whenever $x_0$ belongs to the Lebesgue set of $f$. In case $\varphi_y(t)$ is the Poisson kernel $P(t, y) = c_n[y/(y^2 + |t|^2)^{(n+1)/2}]$ we have the easily established inequality
\begin{equation}\label{eq:c2-3-17}
P(x - t, y) = \varphi_y(x - t) \leq d_\alpha \varphi_y(x_0 - t) = d_\alpha P(x_0 - t, y)
\end{equation}
for $(x, y) \in \Gamma_\alpha(x_0)$, where $d_\alpha^{2/(n+1)} = \max\{1 + 2\alpha^2, 2\}$. Consequently, using Lemma~1.17, part (b) (in Chapter~I),
\[
\begin{aligned}
|u(x, y) - f(x_0)| &= \left|\int_{E_n} f(t) P(x - t, y)\,dt - f(x_0) \int_{E_n} P(x - t, y)\,dt\right| \\
&\leq \int_{E_n} |f(t) - f(x_0)| P(x - t, y)\,dt \leq d_\alpha \int_{E_n} |f(t) - f(x_0)| P(x_0 - t, y)\,dt \\
&= d_\alpha \int_{E_n} |f(t) - f(x_0)| \varphi_y(x_0 - t)\,dt.
\end{aligned}
\]
But, since the last term tends to $0$ as $y$ tends to $0$, the theorem is proved.
\end{proof}

An immediate consequence of the result we have just proved is that the Poisson integral $u$ is nontangentially bounded at almost every point of $E_n$. That is, it is a function $u$ defined in $E_{n+1}^+$ that is bounded in the truncated cones $\Gamma_\alpha(x_0) \cap \{(x, y) \in E_{n+1} : 0 < y \leq 1\}$ for almost all $x_0 \in E_n$. In fact, the maximal function $m_f$ can be used to give the following quantitative expression of this:
\begin{equation}\label{eq:c2-3-18}
\sup_{(x,y) \in \Gamma_\alpha(x_0)} |u(x, y)| \leq d_\alpha m_f(x_0).
\end{equation}
This is a simple consequence of Theorem~\ref{thm:c2-3-10} and inequality \eqref{eq:c2-3-17} involving the Poisson kernel.

The following theorem (a basic tool used in the theory of the $H^p$ spaces developed in Chapter~VI) shows that, for harmonic functions, nontangential limits and nontangential boundedness are essentially equivalent.

\begin{theorem}\label{thm:c2-3-19}
Suppose $u$ is harmonic in $E_{n+1}^+$ and is nontangentially bounded at all the points of a subset $S \subset E_n$ of positive measure, then $u$ has nontangential limits at almost every point of $S$.
\end{theorem}

\begin{proof}
First, we observe that it suffices to show that whenever $S \subset E_n$ is a subset of positive measure at which $u$ is nontangentially bounded then there exists a subset of $S$, also of positive measure, at whose points $u$ has nontangential limits. Considering only cones $\Gamma_\alpha(q)$ with rational apertures $\alpha$ we see that $S$ is a denumerable union of subsets $F = F_{\alpha,m}$ such that $|u(x, y)| \leq m$ for all
\[
(x, y) \in \left\{\bigcup_{q \in F} \Gamma_\alpha(q)\right\} \cap \{p = (x, y) : 0 < y < 2\} = \mathcal{A},
\]
$m = 1, 2, \dots$. That is, $u$ is uniformly bounded in the region $\mathcal{A}$ formed by the union of congruent vertical cones whose vertices are the points of $F$ and are bounded above by the hyperplane $y = 2$.

Now, it is easy to see that it is sufficient for us to show that $u(p)$ tends to a limit when $p$ tends to $q$ within $\Gamma_\alpha(q)$, for almost all $q \in F$. For, if we discard the sets where this does not occur, for each $F = F_{\alpha,m}$ (recall that $\alpha$ is rational and $m$ is a positive integer), we will have discarded a set of measure zero; if $q$ is a point of $S$ not in this set of measure $0$ and $\beta > 0$ any rational aperture, then there exists a positive integer $k$ such that $|u(p)| = |u(x, y)| \leq k$ for $p \in \Gamma_\beta(q)$, $0 < y < 2$; that is, $q \in F_{\beta,k}$. Consequently, $u(p)$ tends to a limit as $p$ tends to $q$ within $\Gamma_\beta(q)$. Thus, a nontangential limit exists for $u$ at all such $q$'s.

A further reduction is obtained by observing that, dividing by $m$, we can assume that $|u| \leq 1$ in $\mathcal{A}$. Moreover, we can also assume that $S$ is contained in a cube of side $1$; since it suffices to prove the theorem for each intersection of $S$ with a member of a mesh of cubes of side $1$.

Since $u$ is harmonic in $E_{n+1}^+$ and bounded in the region $\mathcal{A}$, a standard argument using the mean value property and Harnack's inequality shows that the family $\{u(x, y)\}$ is equicontinuous on compact subsets of $\mathcal{A}$. Applying a diagonal process we can extract a sequence $\{u_j\}$ of such functions, formed by suitable truncations of $u$, which converges to a function $\varphi$ that is the Poisson integral of a function in $L^2(E_n)$. Write $u = \varphi + \psi$, and set $\psi_j = u_j - \varphi_j$. Then the $\psi_j$ satisfy
\begin{enumerate}
\item[(1)] $|\psi_j(p)| \leq 2$ for $p \in \mathcal{D}$;
\item[(2)] $\psi_j(p) \to 0$ as $p \to q \in F$, $p \in \mathcal{D}$.
\end{enumerate}
The theorem will be established if we can find a function $w$, harmonic in $E_{n+1}^+$, satisfying
\begin{enumerate}
\item[(a)] $w(x, y) \geq 0$ in $E_{n+1}^+$;
\item[(b)] $w(x, y) \geq 2$ in $\mathcal{D} - F$;
\item[(c)] $w$ has nontangential limits $0$ at almost every point of $F$.
\end{enumerate}
If this were the case, $w(p) \pm \psi_j(p) \geq 0$ in $\mathcal{D} - F$ because of (1) and (b); and (2) together with (a) gives
\[
\lim_{\substack{p \to q \in F, \\ p \in \mathcal{D}}} \{w(p) \pm \psi_j(p)\} \geq \lim_{\substack{p \to q \in F, \\ p \in \mathcal{D}}} w(p) \geq 0.
\]
It follows that $w(p) \pm \psi_j(p) \geq 0$ for all $p \in \mathcal{D}$, by the minimum principle. Letting $j$ tend to $\infty$ we conclude that $|\psi| \leq w$ in $\mathcal{D}$, and the existence of nontangential limits for $u$ follows.

We shall now construct a function $w$ satisfying the three properties (a), (b), and (c). Let $\xi$ be the characteristic function of $E_n - F$ and let, for $y > 0$,
\[
w(x, y) = 2y + cy \int_{E_n} \frac{\xi(s)}{(y^2 + |x - s|^2)^{(n+1)/2}}\,ds,
\]
where $c$ will be a positive constant (to be chosen later). It is clear that $w \geq 0$ in $E_{n+1}^+$. Property (c) is an immediate consequence of Theorem~\ref{thm:c2-3-16}. Thus, we need only show that $w$ satisfies property (b). We first observe that if $(x, y) \in \mathcal{D} - F$ with $y = 1$ then $w(x, y) \geq 2y = 2$. Now suppose that $p = (x, y) \in \mathcal{D}$ with $0 < y < 1$. Consider the inverted cone with vertex $p$ and aperture $\alpha$. If $C$ is the interior of the sphere formed by the intersection of this cone with the hyperplane $E_n$ then $C$ and $F$ are disjoint. For, if this were not the case, there would exist a point $q \in C \cap F$ and, consequently, $p$ would be in the interior of $\Gamma_\alpha(q)$. But this is impossible since $p \in \mathcal{D}$. Thus,
\[
\begin{aligned}
w(p) &\geq cy \int_C \frac{ds}{(y^2 + |x - s|^2)^{(n+1)/2}} = c\omega_{n-1} y \int_0^{\alpha y} \frac{r^{n-1}\,dr}{(y^2 + r^2)^{(n+1)/2}} \\
&= c\omega_{n-1} \int_0^\alpha \frac{r^{n-1}\,dr}{(1 + r^2)^{(n+1)/2}}.
\end{aligned}
\]
This last expression can obviously be made to exceed $2$ by choosing $c$ appropriately. Hence $w$ satisfies (b) and the theorem is proved.
\end{proof}

We shall be interested in extending this result to functions that are harmonic in sets of variables (these will be called \emph{multiply harmonic} functions). By this we mean twice differentiable functions, $u$, defined on a domain, $\mathcal{D}$, of points $(x^{(1)}, x^{(2)}, \dots, x^{(k)})$ belonging to the Cartesian product $E_{n_1} \times E_{n_2} \times \cdots \times E_{n_k}$ of $k$ Euclidean spaces, that satisfies
\[
\sum_{i=1}^{n_j} \frac{\partial^2 u}{(\partial x_i^{(j)})^2} = 0
\]
for $1 \leq j \leq k$. We shall encounter such functions in the next chapter, where we deal with analytic functions of several variables. In this case, each of the spaces $E_{n_j}$ will be two-dimensional and can be identified with a copy of the complex plane.

It is not hard to state this extension and to indicate the changes in the proof of Theorem~\ref{thm:c2-3-19} that are needed in order to establish it. One of these changes involves results analogous to Theorem~\ref{thm:c2-2-5} and Theorem~\ref{thm:c2-3-16} for ``iterated Poisson integrals'' (i.e., solutions of a variant of the Dirichlet problem involving multiply harmonic functions).

These iterated Poisson integrals are examples of a family of operators somewhat more general than that occurring in Theorem~\ref{thm:c2-3-12}, in that they involve several independent parameters tending to $0$. That is, the family $\{T_\varepsilon\}$ is parametrized by the set $\{\varepsilon = (\varepsilon_1, \varepsilon_2, \dots, \varepsilon_m) \in E_m : \varepsilon_j > 0, j = 1, 2, \dots, m\}$ (the first octant in $E_m$); their behavior as $\varepsilon$ tends to $0 \in E_m$ will be the object of our study. In the case of iterated Poisson integrals, for example, we take
\[
P_j(t, \varepsilon_j) = c_{n_j} \frac{\varepsilon_j}{(|t|^2 + \varepsilon_j^2)^{(n_j+1)/2}}
\]
for $t = (t_1, \dots, t_{n_j}) \in E_{n_j}$, $j = 1, 2$. Then, for $f \in L^p(E_n)$ we define $T_\varepsilon f = T_{(\varepsilon_1, \varepsilon_2)} f$ by letting
\[
(T_\varepsilon f)(x) = \int_{E_{n_1}} P_1(x^{(1)} - t, \varepsilon_1) \left\{\int_{E_{n_2}} P_2(x^{(2)} - s, \varepsilon_2) f(t, s)\,ds\right\} dt,
\]
where $x = (x^{(1)}, x^{(2)}) \in E_{n_1} \times E_{n_2} = E_n$. If $f(x) = f_1(x^{(1)}) f_2(x^{(2)})$ with $f_j \in L^p(E_{n_j}) \cap C_0(E_{n_j})$, $j = 1, 2$, it follows from Theorem~\ref{thm:c2-2-1} that $(T_\varepsilon f)(x)$ converges to $f(x)$ uniformly as $\varepsilon = (\varepsilon_1, \varepsilon_2) \to (0, 0)$. The class of finite linear combinations of such functions is dense in $L^p(E_n)$, $1 \leq p < \infty$; thus, if we can show that inequality \eqref{eq:c2-3-13} holds we can then conclude that $\lim_{\varepsilon \to 0} (T_\varepsilon f)(x)$ exists and is finite.

In order to do this we define $M^{(1)}$ to be the operator of $L^p(E_n)$ assigning to $f$ the maximal function of $f(\,\cdot\,, x^{(2)})$; that is,
\[
(M^{(1)} f)(x) = \sup_{r > 0} \Omega_{n_1}^{-1} r^{-n_1} \int_{|t| \leq r} |f(x^{(1)} - t, x^{(2)})|\,dt.
\]
Similarly, we let
\[
(M^{(2)} f)(x) = \sup_{r > 0} \Omega_{n_2}^{-1} r^{-n_2} \int_{|s| \leq r} |f(x^{(1)}, x^{(2)} - s)|\,ds.
\]
Thus, by Theorem~\ref{thm:c2-3-7},
\begin{equation}\label{eq:c2-3-20}
\|Mf\|_p \leq A\|f\|_p,
\end{equation}
where $A$ depends on $p$, $n_1$ and $n_2$, but is independent of $f \in L^p(E_n)$ (see the discussion preceding \eqref{eq:c2-3-6}). Indeed, by Theorem~\ref{thm:c2-3-7},
\begin{equation}\label{eq:c2-3-21}
\|M^{(2)} M^{(1)} f\|_p \leq b(p, n_2) \|M^{(1)} f\|_p \leq b(p, n_2) b(p, n_1) \|f\|_p = B \|f\|_p.
\end{equation}
On the other hand, it follows from \eqref{eq:c2-3-8} that
\[
\begin{aligned}
|(T_\varepsilon f)(x)| &= \left|(T_{(\varepsilon_1, \varepsilon_2)} f)(x)\right| \\
&\leq \int_{E_{n_2}} P_2(x^{(2)} - s, \varepsilon_2) (M^{(1)} f)(x^{(1)}, s)\,ds \leq (M^{(2)} M^{(1)} f)(x^{(1)}, x^{(2)}).
\end{aligned}
\]
Consequently, $(Mf)(x) = \sup_{\varepsilon_1, \varepsilon_2 > 0} |(T_{(\varepsilon_1, \varepsilon_2)} f)(x)| \leq (M^{(2)} M^{(1)} f)(x)$ and inequality \eqref{eq:c2-3-20} follows from \eqref{eq:c2-3-21}.

We can, therefore, conclude that $\lim (T_\varepsilon f)(x)$ as $\varepsilon = (\varepsilon_1, \varepsilon_2) \to (0, 0)$ exists and is finite for almost all $x \in E_n$. It is clear that we can modify the above argument to apply to $m$-fold iterated Poisson integrals. More precisely, we have the following theorem:

\begin{theorem}\label{thm:c2-3-22}
If $f \in L^p(E_n)$, $1 < p < \infty$, and
\[
\begin{aligned}
&u(x, \varepsilon) = u(x^{(1)}, x^{(2)}, \dots, x^{(m)}; \varepsilon_1, \varepsilon_2, \dots, \varepsilon_m) = \\
&\int_{E_{n_1}} \cdots \int_{E_{n_m}} P_1(x^{(1)} - t^{(1)}, \varepsilon_1) \cdots P_m(x^{(m)} - t^{(m)}, \varepsilon_m) f(t^{(1)}, \dots, t^{(m)})\,dt^{(1)} \cdots dt^{(m)},
\end{aligned}
\]
where $t^{(j)}$ and $x^{(j)}$ belong to $E_{n_j}$ and $\varepsilon = (\varepsilon_1, \varepsilon_2, \dots, \varepsilon_m)$ belongs to the first octant of $E_m$ (that is, $\varepsilon_j > 0$ for $j = 1, 2, \dots, m$), then
\[
\lim_{\varepsilon \to 0} u(x, \varepsilon) = \lim_{\varepsilon_1, \dots, \varepsilon_m \to 0} u(x^{(1)}, x^{(2)}, \dots, x^{(m)}; \varepsilon_1, \dots, \varepsilon_m) = f(x)
\]
for almost every $x \in E_n$.
\end{theorem}

\begin{proof}
We have already seen that $\lim_{\varepsilon \to 0} u(x, \varepsilon)$ exists and is finite almost everywhere. We need only show that this limit is $f(x)$ a.e. As in the proof of Corollary~\ref{cor:c2-3-14} it suffices to show that $\|u(\,\cdot\,, \varepsilon) - f\|_p \to 0$ as $\varepsilon \to 0$. This can be done in much the same way as was done in the proof of Theorem~1.18 in Chapter~I. In order to simplify the notation we shall, again, present the proof when $m = 2$; the method extends in an obvious way to the general case.

For $(h, k) \in E_{n_1} \times E_{n_2}$ let
\[
\omega_{p,f}(h, k) = \left(\int_{E_{n_1}} \int_{E_{n_2}} |f(x^{(1)} + h, x^{(2)} + k) - f(x^{(1)}, x^{(2)})|^p\,dx^{(1)} dx^{(2)}\right)^{1/p}
\]
be the $L^p$ modulus of continuity of $f$. We have seen that $\omega(h, k) \to 0$ as $|h|, |k| \to 0$ and that $\omega(h, k) \leq 2\|f\|_p$. Thus, by Minkowski's inequality for integrals and the Lebesgue dominated convergence theorem,
\[
\left(\int_{E_n} |u(x, \varepsilon) - f(x)|^p\,dx\right)^{1/p} = \int_{E_{n_1}} \int_{E_{n_2}} \omega(-\varepsilon_1 t, -\varepsilon_2 s) P_1(t, 1) P_2(s, 1)\,dt\,ds \to 0
\]
as $\varepsilon_1, \varepsilon_2 \to 0$. This completes the proof.
\end{proof}

A modification of the argument we have just given shows that the theorem (as well as its nontangential version that follows below) is valid also for $p = \infty$. Since, in any case, this is a consequence of Theorem~\ref{thm:c2-3-24} proved below, we need not present this modification.

Theorem~\ref{thm:c2-3-22} asserts that the iterated Poisson integral of a function in $L^p(E_n)$ converges to the function almost everywhere if we approach the points $x_0 = (x_0^{(1)}, x_0^{(2)}, \dots, x_0^{(m)})$ in $E_n$ ``downward'' along the products of the lines $x^{(j)} = x_0^{(j)}$ in $E_{n_j+1}$, $j = 1, 2, \dots, m$. As was the case before, more general limits exist. These are generalizations of nontangential limits: for
\[
\Gamma_{\alpha_j}^{(j)}(x_0^{(j)}) = \{(x^{(j)}, y_j) \in E_{n_j+1}^+ : |x^{(j)} - x_0^{(j)}| < \alpha_j y_j\}, \qquad j = 1, 2, \dots, k,
\]
a function, $u$, defined on $E_{n_1+1}^+ \times \cdots \times E_{n_k+1}^+$, is said to have the nontangential limit $l$ in each set of variables at the point $x_0 = (x_0^{(1)}, x_0^{(2)}, \dots, x_0^{(k)})$ if $u$ tends to $l$ as the point $(x^{(1)}, y_1; x^{(2)}, y_2; \dots; x^{(k)}, y_k)$ tends to $x_0$ within the Cartesian product $\Gamma_{\alpha_1}^{(1)} \times \Gamma_{\alpha_2}^{(2)} \times \cdots \times \Gamma_{\alpha_k}^{(k)}$.

\begin{corollary}\label{cor:c2-3-23}
If $f \in L^p(E_n)$, $1 < p < \infty$, and $u(x, \varepsilon)$ is the function defined in Theorem~\ref{thm:c2-3-22}, then $u(x, \varepsilon)$ has the nontangential limit $f(x_0)$ in each set of variables at almost all $x_0 \in E_n$.
\end{corollary}

We say that $u$ is \emph{nontangentially bounded in each set of variables} at $x_0 \in E_n$ if $u$ is bounded in the Cartesian product of the truncated cones $\Gamma_{\alpha_j}(x_0^{(j)}) \cap \{(x^{(j)}, y_j) : 0 < y_j \leq 1\}$, $1 \leq j \leq k$, for all $k$-tuples $(\alpha_1, \dots, \alpha_k)$ of positive real numbers. In complete analogy to the case of harmonic functions on $E_{n+1}^+$, this nontangential boundedness is almost everywhere equivalent to the existence of nontangential limits in each set of variables, even for harmonic functions that are not Poisson integrals. More precisely, we have the following extension of Theorem~\ref{thm:c2-3-19}:

\begin{theorem}\label{thm:c2-3-24}
Suppose $n_1 + n_2 + \cdots + n_k = n$ and $u$ is a function defined in the Cartesian product $E_{n_1+1}^+ \times E_{n_2+1}^+ \times \cdots \times E_{n_k+1}^+$ that is harmonic in each set of variables $(x^{(j)}, y_j) = (x_1^{(j)}, \dots, x_{n_j}^{(j)}, y_j)$ in $E_{n_j+1}^+$, $1 \leq j \leq k$. (That is, $u$ is twice differentiable and satisfies
\[
\Delta_j u = \sum_{i=1}^{n_j} \frac{\partial^2 u}{(\partial x_i^{(j)})^2} + \frac{\partial^2 u}{\partial y_j^2} = 0.)
\]
If $u$ is nontangentially bounded in each set of variables at all the points of a subset $S \subset E_n$ of positive measure, then $u$ has nontangential limits in each set of variables at almost every point of $S$.
\end{theorem}

\begin{proof}
The main ideas in establishing this theorem were developed in the proof of Theorem~\ref{thm:c2-3-19}. Therefore, we shall only indicate the changes in that proof that are needed to obtain the present theorem. We confine ourselves to the case $k = 2$ (the argument extends easily to the general case).

In analogy with the notation used in the proof of Theorem~\ref{thm:c2-3-19} we let $p = (p_1, p_2) = (x^{(1)}, y_1; x^{(2)}, y_2)$ denote the general point of $E_{n+1}^+ \times E_{n+1}^+$; and let $q = (q_1, q_2) = (x^{(1)}, 0; x^{(2)}, 0)$ denote the general point of $E_n \times E_n = E_{2n}$ (the latter is the distinguished boundary of $E_{n+1}^+ \times E_{n+1}^+$). By reasoning similar to that used before, we essentially reduce the problem to showing that if $|u| \leq 1$ in the region
\[
\mathcal{D} = \left\{\bigcup_{(q_1,q_2) \in F} \gamma_\alpha(q_1, q_2)\right\} \cap \{p = (x^{(1)}, y_1; x^{(2)}, y_2) : 0 < y_1, y_2 < 1\},
\]
then $u(p)$ tends to a limit when $p$ tends to $q = (q_1, q_2)$ within $\gamma_\alpha(q) = \gamma_\alpha(q_1, q_2)$, for almost all $q \in F \subset E_{2n}$. Here we have written $\gamma_\alpha(q_1, q_2)$ for the Cartesian product cone $\Gamma_{\alpha_1}(q_1) \times \Gamma_{\alpha_2}(q_2)$. By taking appropriate subsets we can assume that $F$ is closed and, in addition, that $F$ is contained in the unit sphere of $E_{2n}$.

We let $\xi$ be the characteristic function of that part of $E_{2n} - F$ contained in a fixed sphere of sufficiently large radius centered at the origin. We then define $w = 2(y_1 + y_2) + cW(x^{(1)}, y_1; x^{(2)}, y_2)$, where $W$ is the repeated Poisson integral of $\xi$, and $c$ is a positive constant to be determined later. Condition (a) is obvious, and condition (c) -- the vanishing of the nontangential limits on almost every point of $F$ -- follows from Corollary~\ref{cor:c2-3-23} (since $\xi$ belongs to all the spaces $L^p(E_{2n})$, $1 < p < \infty$). Thus, we need only prove (b), and it suffices to consider the case in which $p = (x^{(1)}, y_1; x^{(2)}, y_2)$ satisfies $0 < y_1 < 1$, $0 < y_2 < 1$. Consider the Cartesian product of the inverted cones with vertices $p_1$ and $p_2$ and apertures $\alpha_1$ and $\alpha_2$. These cones intersect the hyperplanes $y_1 = 0$ and $y_2 = 0$ in two open spheres $C_1$ and $C_2$. The product $C_1 \times C_2$ does not contain a point of $F$ (if this were not the case $p$ would belong to $\mathcal{D}$), and, since the spheres $C_1$ and $C_2$ are centered at $x^{(1)}$ and $x^{(2)}$ and have radii $\alpha_1 y_1$ and $\alpha_2 y_2$ respectively,
\[
\begin{aligned}
w(p) &\geq cW(p) \geq c \prod_{j=1}^2 \left(c_n y_j \int_{C_j} \frac{ds^{(j)}}{(y_j^2 + |x^{(j)} - s^{(j)}|^2)^{(n+1)/2}}\right) \\
&= c \prod_{j=1}^2 \left(y_j c_n \omega_{n-1} \int_0^{\alpha_j y_j} \frac{r^{n-1}\,dr}{(y_j^2 + r^2)^{(n+1)/2}}\right) \geq 2,
\end{aligned}
\]
if $c$ is large enough. The remaining boundary points, where one of the $y_j$ equals $0$, are handled by Fatou's lemma together with the estimate
\begin{equation}\label{eq:c2-3-25}
\liminf_{p_2' \to p_2} c_n \int_{E_n} \frac{y_2'}{|p_2' - q_2|^{(n+1)}} \xi(q_1, q_2)\,dq_2 \geq \xi(q_1, p_2),
\end{equation}
which holds because $\xi$ is continuous at every point of the open set where it equals $1$. Applying the argument of Theorem~\ref{thm:c2-3-19} to the functions $\psi_j$ in the region $\mathcal{D}_{p_2}$ then completes the proof of (b), and with it the majorization
\begin{equation}\label{eq:c2-3-26}
|\psi(p)| \leq w(p), \qquad p \in \mathcal{D}.
\end{equation}
To conclude the proof of the theorem, we shall need to extend \eqref{eq:c2-3-26} so that for almost all $q \in F$, $|\psi(p)| \leq w(p)$ where $p$ may lie in product cones $\Gamma_{\beta_1}(q_1) \times \Gamma_{\beta_2}(q_2)$, with arbitrary apertures $\beta_1$ and $\beta_2$, as long as $p$ is sufficiently close to $q = (q_1, q_2)$. This will indeed be the case at all points of strong density of $F$. These are the points which satisfy the conclusion of the following lemma:
\end{proof}

\begin{lemma}\label{lem:c2-3-27}
Let $\xi$ be the characteristic function of an arbitrary measurable set $F$ of $E_{2n} = E_n \times E_n$. Then for almost every $x = (x^{(1)}, x^{(2)}) \in F$
\[
\lim_{\substack{h_1 \to 0 \\ h_2 \to 0}} \Omega_n^{-2} h_1^{-n} h_2^{-n} \int_{|s^{(1)}| < h_1, |s^{(2)}| < h_2} \int \xi(x^{(1)} - s^{(1)}, x^{(2)} - s^{(2)})\,ds^{(1)} ds^{(2)} = 1.
\]
\end{lemma}

\begin{proof}
For any $f \in L^p(E_n \times E_n)$, we define the \emph{strong maximal function} $\mu_f$ by
\[
\mu_f(x^{(1)}, x^{(2)}) = \sup_{h_1 > 0, h_2 > 0} \Omega_n^{-2} h_1^{-n} h_2^{-n} \int_{|s^{(1)}| < h_1, |s^{(2)}| < h_2} \int |f(x^{(1)} - s^{(1)}, x^{(2)} - s^{(2)})|\,ds^{(1)} ds^{(2)}.
\]
Then, as in the proof of \eqref{eq:c2-3-20}, $\mu_f(x) \leq (M^{(2)} M^{(1)} f)(x)$, and so
\begin{equation}\label{eq:c2-3-28}
\|\mu_f\|_p \leq A_p \|f\|_p, \qquad 1 < p < \infty.
\end{equation}
By the same argument as in the proof of Theorem~\ref{thm:c2-3-22} it then follows that
\begin{equation}\label{eq:c2-3-29}
\lim_{\substack{h_1 \to 0 \\ h_2 \to 0}} \Omega_n^{-2} h_1^{-n} h_2^{-n} \int_{|s^{(1)}| < h_1, |s^{(2)}| < h_2} \int |f(x^{(1)} - s^{(1)}, x^{(2)} - s^{(2)}) - f(x^{(1)}, x^{(2)})|\,ds^{(1)} ds^{(2)} = 0
\end{equation}
for almost every $x \in E_n \times E_n$ whenever $f \in L^p(E_n \times E_n)$, $1 < p < \infty$. Taking $f = \xi$ completes the proof of the lemma.
\end{proof}

\section{Subharmonic Functions and Majorization by Harmonic Functions}

A function $s : \mathcal{D} \to [-\infty, \infty)$ defined on a domain $\mathcal{D} \subset E_n$ is said to be \emph{subharmonic} if:
\begin{enumerate}
\item[(1)] $s$ is upper semicontinuous on $\mathcal{D}$;
\item[(2)] $s$ is not identically $-\infty$ on any component of $\mathcal{D}$;
\item[(3)] For every $x \in \mathcal{D}$ and every $r > 0$ such that $\overline{B_r(x)} \subset \mathcal{D}$,
\begin{equation}\label{eq:c2-4-1}
s(x) \leq \mathcal{M}_{x,s}(r) = \frac{1}{\omega_{n-1}} \int_\Sigma s(x + rt')\,dt'.
\end{equation}
\end{enumerate}
The basic facts concerning subharmonic functions are contained in the following two theorems.

\begin{equation}\label{eq:c2-4-2}
\Delta s \geq 0 \quad \text{throughout } \mathcal{D}.
\end{equation}

\begin{theorem}\label{thm:c2-4-3}
If $s \in C^2(\mathcal{D})$, then $s$ is subharmonic in $\mathcal{D}$ if and only if $\Delta s \geq 0$ throughout $\mathcal{D}$.
\end{theorem}

\begin{theorem}\label{thm:c2-4-4}
Suppose $s \geq 0$ is continuous in a domain $\mathcal{D}$ and satisfies the mean value inequality $s(x) \leq \mathcal{M}_{x,s}(r)$ whenever the sphere of radius $r$ about $x$ is contained in $\mathcal{D}$. Then $s$ is subharmonic in $\mathcal{D}$.
\end{theorem}

\begin{customthm}{4.4$'$}\label{thm:c2-4-4p}
Suppose $s \geq 0$ is continuous in a domain $\mathcal{D}$ and is subharmonic. If $\Phi : [0, \infty) \to [0, \infty)$ is a non-decreasing convex function, then $\Phi(s)$ is subharmonic in $\mathcal{D}$. In particular, $s^p$ is subharmonic for all $p \geq 1$.
\end{customthm}

\begin{theorem}\label{thm:c2-4-5}
Let $F = (u_1, u_2, \dots, u_k)$ be a vector of harmonic functions in $\mathcal{D}$ satisfying the generalized Cauchy--Riemann equations of Chapter~VI. Then $|F|^p = (\sum_{j=1}^k u_j^2)^{p/2}$ is subharmonic in $\mathcal{D}$ for all $p \geq (n-1)/n$.
\end{theorem}

\begin{theorem}[Least Harmonic Majorant]\label{thm:c2-4-6}
Let $s$ be subharmonic in a domain $\mathcal{D}$. If there exists a harmonic function $H$ on $\mathcal{D}$ such that $s(x) \leq H(x)$ for all $x \in \mathcal{D}$, then there exists a unique \emph{least harmonic majorant} $h$ of $s$; that is, $h$ is harmonic in $\mathcal{D}$, $s \leq h$, and if $H$ is any other harmonic majorant of $s$ then $h \leq H$.
\end{theorem}

\begin{equation}\label{eq:c2-4-7}
h(x) = \lim_{r \to 1} \frac{1}{\omega_{n-1}} \int_S P(x/r, t') s(rt')\,dt'.
\end{equation}

\begin{equation}\label{eq:c2-4-8}
s(x, y) \leq \int_{E_n} P(x - t, y) f(t)\,dt.
\end{equation}

\begin{equation}\label{eq:c2-4-9}
\lim_{y \to 0} \int_{E_n} s(x, y)^p\,dx < \infty.
\end{equation}

\begin{equation}\label{eq:c2-4-10}
\sup_{y > 0} \int_{E_n} |u(x, y)|^p\,dx < \infty.
\end{equation}

\begin{equation}\label{eq:c2-4-11}
u(x, y) = \int_{E_n} P(x - t, y) f(t)\,dt + c y.
\end{equation}

\begin{equation}\label{eq:c2-4-12}
\sup_{(t,y) \in \Gamma_\alpha(x)} |F(t, y)| \leq C_\alpha M(|f|)(x).
\end{equation}

A particularly important class of subharmonic functions is provided by the following: if $F = (u_0, u_1, \dots, u_n)$ is a system of conjugate harmonic functions on $E_{n+1}^+$, then $|F|^p$ is subharmonic for $p \geq (n-1)/n$; this will be established in Chapter~VI and will be used there in order to develop the theory of the spaces $\mathbf{H}^p$ of conjugate harmonic functions.

\section{Further Results}

\paragraph{5.1.} At the beginning of this chapter we introduced the mean value $\mathcal{M}_{x,u}(r) = \mathcal{M}(r)$ of a function $u$ taken over the surface of the sphere $S_r(x)$ of radius $r$ centered at $x$. It is just as natural to form the mean value
\[
\mathcal{A}_{x,u}(r) = \mathcal{A}(r) = \frac{1}{\Omega_n r^n} \int_{|t| \leq r} u(x + t)\,dt,
\]
of the function $u$, taken over the volume of the sphere $S_r(x)$. We can, therefore, consider two different mean-value properties: the one introduced at the beginning of this chapter and the analogous one involving $\mathcal{A}(r)$. The argument used at the beginning of this chapter shows that they are equivalent.

\paragraph{5.2.} There are several results concerning harmonic (and analytic) functions defined on unbounded domains that are extensions of the maximum principle as stated in Corollary~\ref{cor:c2-1-3p}. These results can be described in a general way as follows: Suppose we are given a harmonic function $u$ in $E_{n+1}^+$ satisfying the growth conditions $u(x, y) = o(e^{a|x|})$ and $u(x, y) = o(y)$ as $|x| \to \infty$ and $y \to \infty$; then
\begin{enumerate}
\item[(a)] $u(x, 0) \leq A < \infty$ for all $x \in E_n$ implies $u(x, y) \leq A$ for all $(x, y) \in E_{n+1}^+$;
\item[(b)] $u(x, 0) \geq B > -\infty$ for all $x \in E_n$ implies $u(x, y) \geq B$ for all $(x, y) \in E_{n+1}^+$;
\item[(c)] $|u(x, 0)| \leq C < \infty$ for all $x \in E_n$ implies $|u(x, y)| \leq C$ for all $(x, y) \in E_{n+1}^+$.
\end{enumerate}
To show, say, part (b) we can assume $B = 0$ since the general case can be reduced to this one by considering the function $u(x, y) - B$. Fix $\delta > 0$ and let $g(x, y) = \sin(a\sqrt{n}(y + \delta)) \prod_{k=1}^n \cosh a x_k$; then $g$ is a harmonic function which is positive on the set where $y \geq 0$ and, since
\[
\prod_{k=1}^n \cosh a x_k \geq 2^{-n} \prod_{k=1}^n e^{a|x_k|} = 2^{-n} \exp\left\{a\sum_{k=1}^n |x_k|\right\} \geq 2^{-n} \exp\{a|x|\},
\]
it grows exponentially at infinity. These estimates, together with the hypotheses $u(x, y) = o(e^{a|x|})$ and $u(x, y) = o(y)$, imply the desired inequality $u(x, y) + \varepsilon y + \eta g(x, y) \geq 0$ for $|x| \geq N$ (if $N$ is sufficiently large) and $0 \leq y \leq y_0$. By the minimum principle we obtain the result.

\paragraph{5.3.} We showed (see Lemma~\ref{lem:c2-2-6} and inequality \eqref{eq:c2-4-8}) that if a (nonnegative) subharmonic function $s$ in $E_{n+1}^+$ satisfies $\|s(\,\cdot\,, y)\|_p \leq c$ for all $y > 0$ and some $p \geq 1$, then $\|s(\,\cdot\,, y)\|_\infty = O(y^{-n/p})$. As a consequence of this, if $q > p$ then $s(\,\cdot\,, y) \in L^q(E_n)$ for all $y > 0$ and $\|s(\,\cdot\,, y)\|_q = O(y^{-n(q-p)/(pq)})$. Indeed,
\[
\int_{E_n} [s(x, y)]^q\,dx = \int_{E_n} [s(x, y)]^p [s(x, y)]^{q-p}\,dx \leq \|s(\,\cdot\,, y)\|_\infty^{q-p} c^p \leq A[y^{-n/p}]^{q-p},
\]
and the desired inequality follows by taking $q$-th roots.

\paragraph{5.4.} Theorems~\ref{thm:c2-3-19} and~\ref{thm:c2-3-24} indicate that the concept of nontangential convergence at the boundary is essential. That this is indeed the fact can be seen from the following two results dealing with functions defined on the unit disk.

\begin{enumerate}
\item[(i)] Let $C_\theta$ be a simple closed curve passing through the point $z = 1$, all of whose other points are contained in the interior of the unit circle $K = \{z = x + iy : |z| = 1\}$. Suppose, further, that $C_\theta$ is tangent to the unit circle at $1$. Then there exists an analytic function $F$, defined in the interior of $K$, such that, for almost all $\theta$, $F(z)$ does not tend to any limit as $z$ tends to $e^{i\theta}$ within $C_\theta$ (see Zygmund~\cite{Zygmund1}, Vol.~I, page~280).
\item[(ii)] Let $G$ be a continuous function defined in the interior of the unit circle, $K$, and let $E$ be a set of the first category in $K$. Then there exists an analytic function $F$ defined in the interior of $K$ such that
\[
\lim_{r \to 1} \{F(re^{i\theta}) - G(re^{i\theta})\} = 0
\]
whenever $e^{i\theta} \in E$ (see Zygmund~\cite{Zygmund1}, Vol.~II, page~204).
\end{enumerate}
Since $E$ can have measure $2\pi$, this last result asserts that the radial behavior, at almost all points of $K$, of analytic functions is no better than that of continuous functions. Thus, for example, an analytic function $F$ can tend to 0 along almost every radius and yet be unbounded in every sector.

\paragraph{5.5.} Theorem~\ref{thm:c2-3-19} has several extensions. Carleson (see \cite{Carleson1}) has shown that the same result holds if we assume that the (real-valued) harmonic function $u$ is nontangentially bounded from below or from above at each point of $S$. In the two-dimensional case this is equivalent to a classical result of Plessner (see Theorem~\ref{thm:c2-3-24}).

\paragraph{5.6.} By slightly modifying the argument used in the proof of Lemma~\ref{lem:c2-3-3} we can obtain the following closely related result: Suppose that to each point $x$ of the bounded set $F \subset E_n$ there corresponds a nonempty open sphere $S_x$ with center $x$. Then there exists a sequence $\{x_j\} \subset F$ such that $F \subset \bigcup_j S_{x_j}$ and each point of $E_n$ belongs to at most $a_n$ of the spheres $S_{x_j}$, where $a_n$ depends only on $n$. If $\nu_1$ is a measure, $\nu_1 \ll$ Lebesgue measure and $\nu_1(S_x) = \int_{S_x} f$ for some integrable $f$, then $\varphi(x) = \sup_x \nu_1(S_x)/|S_x|$ is the Hardy--Littlewood maximal function $m_f$. Then, if $F_s = \{x \in E_n : \varphi(x) > s > 0\}$, we have $\nu_2(F_s) \leq a_n \nu_1(E_n) s^{-1}$.

\paragraph{5.7.} In Theorem~\ref{thm:c2-3-22} and Corollary~\ref{cor:c2-3-23} we assumed $1 < p < \infty$. The case $p = \infty$ was discussed at the end of the proof of Theorem~\ref{thm:c2-3-22}. There, we pointed out that the result in this case follows from similar arguments. The case $p = 1$ is quite different, however: the iterated Poisson integrals
\[
u(x_1, \varepsilon_1, x_2, \varepsilon_2) = \int_{-\infty}^\infty \int_{-\infty}^\infty f(x_1 - t_1, x_2 - t_2) P(t_1, \varepsilon_1) P(t_2, \varepsilon_2)\,dt_1 dt_2
\]
do not converge to $f(x_1, x_2)$ as $\varepsilon_1, \varepsilon_2 \to 0$ for almost all $x = (x_1, x_2) \in E_2$ (see Jessen, Marcinkiewicz, and Zygmund~\cite{JessenMarcinkiewiczZygmund1}, where the iterated Poisson integral for the unit circle is discussed. The argument used there can be adapted to the present situation).

\paragraph{5.8.} It follows easily from our development of subharmonic functions that if $s$ has continuous second partial derivatives then the following three conditions are equivalent: (i) $s$ is subharmonic; (ii) $\Delta s \geq 0$; (iii) if $\mathcal{R}$ is a bounded domain with $\overline{\mathcal{R}} \subset \mathcal{D}$ and $h$ is the harmonic function in $\mathcal{R}$ with boundary values $s$ on $\partial\mathcal{R}$ (given by the solution of the Dirichlet problem), then $s \leq h$ in $\mathcal{R}$. The extension of these equivalences to the case in which $s$ is merely upper semicontinuous is obtained by an approximation argument.

\bibnoteshead

Kellog~\cite{Kellog1} is a comprehensive text on the theory of harmonic functions of three variables. For a more modern treatment of the general $n$-dimensional case see Brelot~\cite{Brelot1}. The characterization of Poisson integrals is a classical result and the proofs used here are not very different from those used for the characterization of the Ces\'aro and Abel means of trigonometric series that can be found in Zygmund~\cite{Zygmund1}, Chapter~III. Theorem~3.10 and inequality (3.18) are an extension to upper half-spaces of the corresponding results for Poisson integrals of functions defined on the unit circle originally obtained by Hardy and Littlewood~\cite{HardyLittlewood1}; for these and other variants see K.~T.~Smith [1].% TODO: the OCR'd bibliography lacks a Smith entry; original prints "[1]".
Theorems 3.19 and 3.24 are due to Calder\'on~\cite{Calderon1}. The maximal function was first introduced, in one dimension, by Hardy and Littlewood~\cite{HardyLittlewood1}. The extension to several variables is due to Wiener~\cite{Wiener2}. In connection with the covering lemma used here see Wiener~\cite{Wiener2}, Marcinkiewicz and Zygmund~\cite{MarcinkiewiczZygmund1}, Besicovitch~\cite{Besicovitch2}, Cotlar~\cite{Cotlar1}, and de~Guzm\'an~\cite{deGuzman1}. Theorem~3.12 is closely related to a result of Cotlar~\cite{Cotlar1}; an earlier and more general result is due to Banach~\cite{Banach1}. For a comprehensive treatment of the theory of subharmonic functions of two variables see Rado~\cite{Rado1}. Theorem~4.6 can be found in Stein and G.~Weiss~\cite{SteinWeiss1}. For a more complete treatment of strong differentiability of integrals see Saks~\cite{Saks1}.
